\subsection{全体概要：この章で学ぶこと}
この章は、ソフトウェア開発におけるセキュリティの考え方を扱う。セキュリティは、正しさや信頼性とは別に後で確認する付録ではない。攻撃者、誤用、悪意ある入力、第三者部品、運用中の変更を前提に、開発ライフサイクル全体で扱うべき品質特性である。

到達目標は、ソフトウェアセキュリティ、情報セキュリティ、サイバーセキュリティの違いを説明できること、セキュリティ要求・設計・実装・テスト・脆弱性管理の流れを説明できること、クラウド、IoT、機械学習のような領域固有の注意点を説明できることである。

\begin{itemize}[leftmargin=2em]
\item セキュリティは、要求段階で何を守るかを決め、設計段階で守り方を決め、実装段階で危険な書き方を避け、テスト段階で脆弱性を探し、保守段階で新しい脆弱性に対応する活動である。
\item 脆弱性は、攻撃者が悪用できる誤りである。ソフトウェア本体だけでなく、ライブラリ、市販既製ソフトウェア（COTS）、オペレーティングシステムにも入り得る。
\end{itemize}
\par\smallskip
\begin{center}
\centering
\IfFileExists{assets/generated_figures/ch13/fig-ch13-01.png}{\includegraphics[width=0.96\linewidth]{assets/generated_figures/ch13/fig-ch13-01.png}}{\fbox{\parbox[c][48mm][c]{0.9\linewidth}{\centering 画像生成待ち\\13 章の全体地図}}}
\par\smallskip
\refstepcounter{figure}\label{fig:ch13-01}図 \thefigure: 13 章の全体地図
\end{center}
図 \thefigure は、13 章の全体地図を視覚的に整理したものである。
\par\smallskip
